\begin{frame}<1-4>[t]{学习目标}
\begin{campuscanvas}
% Source slide 2, objects 2; visible 2-4
\only<2-4|handout:1>{\node[anchor=north west,inner sep=0,text=black] at (70.000,200.000) {\begin{minipage}[t]{142.500mm}\raggedright\fontsize{14}{21.00}\selectfont \textcolor{campusgreen}{\textbf{01}}\quad 理解集合的含义，能判断给定对象能否构成集合。\end{minipage}};}
% Source slide 2, objects 15; visible 3-4
\only<3-4|handout:1>{\node[anchor=north west,inner sep=0,text=black] at (70.000,340.000) {\begin{minipage}[t]{142.500mm}\raggedright\fontsize{14}{21.00}\selectfont \textcolor{campusgreen}{\textbf{02}}\quad 掌握集合的常用表示方法，能根据集合特点选择合适方法表示集合。\end{minipage}};}
% Source slide 2, objects 18; visible 4
\only<4|handout:1>{\node[anchor=north west,inner sep=0,text=black] at (70.000,490.000) {\begin{minipage}[t]{142.500mm}\raggedright\fontsize{14}{21.00}\selectfont \textcolor{campusgreen}{\textbf{03}}\quad 熟知常见数集的记法，会用区间表示集合。\end{minipage}};}
\end{campuscanvas}
\end{frame}
